\documentclass[11pt,a4paper]{article}

\usepackage[utf8]{inputenc}
\usepackage[T1]{fontenc}
\usepackage{lmodern}
\usepackage{geometry}
\usepackage{amsmath,amssymb}
\usepackage{graphicx}
\usepackage{booktabs}
\usepackage[hidelinks]{hyperref}
\usepackage{enumitem}
\usepackage{setspace}

% Any other packages
\usepackage{float}

\geometry{margin=2.5cm}
\onehalfspacing

\title{
Link Prediction in Citation Networks using Graph Neural Networks\\
\large Master-Level Graph Neural Networks Course Project Proposal
}

\author{
S\'andor Baranyi \\ \texttt{ct9xfj@inf.elte.hu}
\and
Bal\'azs Bord\'acs \\ \texttt{adjdme@inf.elte.hu}
\and
M\'at\'e Kov\'acs \\ \texttt{u5bky4@inf.elte.hu}
\and
Adorj\'an Nagy-Mohos \\ \texttt{d5vd5e@inf.elte.hu}
}

\date{\today}

\begin{document}

\maketitle

\begin{abstract}
This project studies link prediction in citation networks as a proxy task for citation recommendation.
Using benchmark datasets such as Cora and PubMed, we compare classical link prediction methods
with simple Graph Neural Network (GNN) baselines under a fully reproducible, CPU-friendly setup.
The project evaluates whether feature-aware message-passing models, specifically GCN and GraphSAGE,
improve over non-neural approaches such as local link heuristics and node2vec-based edge scoring.
Our primary research lever is the effect of citation directionality, comparing directed and symmetrized
graph constructions. As a secondary lever, we investigate whether lightweight structural or positional
augmentations improve prediction quality. The outcome of the project is not architectural novelty,
but a careful and reproducible comparison of which modeling choices matter most for citation link prediction.
\end{abstract}

\section{Introduction}

Graph-structured data appear in many real-world applications,
including social networks, biological networks, and citation networks.
Citation graphs represent relationships between scientific papers,
where nodes correspond to papers, and edges represent citation links.

Predicting missing links in citation networks can be interpreted as
recommending relevant references for a given paper.
This task is important for recommendation systems, knowledge graph completion,
and scientific information retrieval.

Graph Neural Networks (GNNs) have recently shown     strong performance on graph learning tasks.
In this project, we investigate whether simple GNN architectures
can outperform classical link prediction methods
on citation graph benchmarks.

\section{Problem Definition}

We model a citation network as a graph $G=(V,E)$, where nodes correspond to papers
and directed edges correspond to citation relations.
The project studies \emph{link prediction}: given a pair of papers $(u,v)$,
the model predicts whether there should be a citation link from $u$ to $v$.

In our setting, this serves as a controlled proxy for citation recommendation.
More precisely, the task is to recover held-out citation edges among papers already present
in the graph, rather than to build a production-grade scholarly recommender for newly written papers.
This framing keeps the problem well defined, benchmarkable, and feasible within the compute budget
of the course.

Formally, let $E_{\mathrm{train}} \subset E$ denote the observed training edges.
The goal is to learn a scoring function
\[
s(u,v) \in \mathbb{R}
\]
that assigns higher scores to true but hidden citation edges than to non-edges.
Positive examples are sampled from held-out citation edges, while negative examples are sampled
from node pairs without observed citation links.

\section{Objectives and Research Questions}

The main objective of the project is to determine whether simple GNN-based link prediction models
provide measurable benefits over classical methods on citation network benchmarks.

\subsection{Research Questions}
\begin{itemize}
    \item RQ1: Do feature-aware GNN models outperform classical link prediction baselines on citation graphs?
    \item RQ2: Does preserving citation directionality improve predictive performance compared to symmetrized graphs?
    \item RQ3: Do lightweight structural or positional augmentations improve the performance of simple GNN baselines?
    \item RQ4: Can all experiments be executed reproducibly on CPU within the project constraints?
\end{itemize}

\subsection{Hypotheses}
\begin{itemize}
    \item H1: GCN and GraphSAGE will outperform local heuristic baselines because they combine graph structure with node features.
    \item H2: Directed modeling will outperform undirected modeling because citation relations are inherently asymmetric.
    \item H3: Adding lightweight structural or positional information will improve link prediction over plain message passing.
    \item H4: A controlled experimental pipeline based on Cora and PubMed, together with lightweight baselines and shallow GNN models, will be feasible to run reproducibly on CPU without requiring specialized hardware.
\end{itemize}

\section{Novelty and Motivation}

Citation datasets are often used for node classification, whereas our project emphasizes link prediction instead.

This project focuses on:
\begin{itemize}
\item link prediction instead of node classification,
\item comparison of classical and neural methods,
\item structural and positional encodings,
\item directed vs. undirected citation graphs.
\end{itemize}

The goal is to analyze how much benefit GNNs provide
over simpler graph algorithms.

\section{Dataset Audit}

We use two standard citation-network benchmarks from the Planetoid family: Cora and PubMed, which are well suited to this project because they provide node features, are widely used in graph learning,
and remain small enough for reproducible CPU-based experiments.

In the default PyTorch Geometric Planetoid representation, citation edges are stored in a symmetrized form,
so the raw \texttt{edge\_index} contains both directions of an undirected connection.
For directionality experiments, we plan to reconstruct a directed variant from the original citation relations, subject to data availability in the benchmark source.

\begin{table}[h]
\centering
\begin{tabular}{lccccccc}
\toprule
Dataset & Nodes & Edge entries & Unique links & Features & Classes & Sparsity & Avg.\ degree \\
\midrule
Cora   & 2708  & 10556 & 5278  & 1433 & 7 & 98.73\% & 3.90 \\
PubMed & 19717 & 88648 & 44324 & 500  & 3 & 89.98\% & 4.50 \\
\bottomrule
\end{tabular}
\caption{Summary of the datasets used in our experiments.}
\end{table}

Cora uses sparse binary bag-of-words features, while PubMed provides denser TF--IDF style features.
Both graphs are sparse, but PubMed is structurally simpler: it forms a single connected component,
whereas Cora contains multiple components.

These datasets are appropriate for controlled comparison between classical and GNN-based link prediction methods,
but they remain simplified benchmarks rather than realistic large-scale citation recommendation systems.
Since the standard Planetoid masks are designed for node classification, our project will use a fixed edge-level split for link prediction.

\section{Literature Review}

Graph Neural Networks have become a central approach for learning on relational data,
because they combine node attributes with graph topology through iterative neighborhood aggregation.
For citation networks in particular, this is appealing because both text-derived node features
and citation structure are informative.
Early benchmark work on citation graphs established datasets such as Cora, CiteSeer, and PubMed
as standard testbeds for graph learning, and showed that these networks exhibit strong homophily
and meaningful structural regularities.
In particular, prior analysis reported high assortativity for Cora and PubMed,
which supports the intuition that neighboring papers often belong to related topics.
This makes citation networks a natural setting for feature-aware graph methods
\cite{sen2008collective,namata2012query}.

A major milestone in this literature is the Graph Convolutional Network (GCN) of
Kipf and Welling \cite{kipf2017semi}.
GCN simplified spectral graph convolutions into a practical message-passing model
that efficiently aggregates information from local neighborhoods.
Their work demonstrated strong node classification performance on citation datasets such as
Cora and PubMed, showing that even shallow neural architectures can exploit both graph structure
and node features effectively.
Although the original GCN paper focused on semi-supervised node classification rather than link prediction,
its importance for our project is twofold.
First, it provides a simple and widely accepted GNN baseline.
Second, its layer design captures the basic intuition that each node representation should be refined
by mixing information from adjacent nodes.
Because our project aims to remain CPU-friendly and reproducible,
GCN is especially attractive as a baseline architecture:
it is conceptually simple, computationally lightweight, and well supported in modern libraries.

A closely related but conceptually different line of work is GraphSAGE,
introduced by Hamilton, Ying, and Leskovec \cite{hamilton2017inductive}.
Whereas GCN is commonly discussed in a transductive setting,
GraphSAGE was designed explicitly for inductive learning:
instead of learning an embedding per node,
it learns aggregation functions that can generate embeddings for previously unseen nodes.
This distinction is relevant for recommendation-style applications,
although our project itself focuses on held-out links among papers already present in the graph.
Even in this transductive setting, GraphSAGE remains a useful contrasting aggregation baseline.
GraphSAGE also broadens the design space by allowing different aggregation operators,
including mean, pooling, and LSTM-based variants.
The original paper showed that GraphSAGE outperformed strong baselines on evolving citation data,
and emphasized that node features and local neighborhood patterns can be jointly exploited.
For our project, GraphSAGE is therefore a natural second GNN baseline:
it remains simple enough for course-level experiments,
while broadening the comparison with a conceptually different neighborhood aggregation strategy.

Before the rise of GNNs, graph representation learning for link prediction
was dominated by unsupervised embedding methods and handcrafted structural heuristics.
Among these, node2vec \cite{grover2016node2vec} is one of the most influential.
node2vec learns continuous node embeddings by optimizing neighborhood co-occurrence
under biased random walks, with parameters controlling a trade-off between breadth-first
and depth-first exploration.
This flexibility allows the method to emphasize either structural equivalence
or community-based similarity, depending on the graph.
Importantly for our project, node2vec was explicitly evaluated on link prediction tasks,
where edge representations are formed by combining pairs of node embeddings.
Because it is scalable, conceptually elegant, and remains competitive in many settings,
node2vec is an essential classical baseline.
It also provides a useful contrast with GNNs:
where node2vec learns embeddings from graph structure alone,
GCN and GraphSAGE can additionally exploit the provided bag-of-words or TF--IDF features.
This comparison directly matches one of our central research questions,
namely whether feature-aware neural message passing yields clear gains
over strong structure-driven baselines.

A different classical perspective comes from graph heuristics and ranking methods.
PageRank \cite{page1999pagerank}, although originally developed to rank web pages,
is also fundamental in citation analysis because it measures node importance
through recursive influence propagation.
While PageRank is not a link predictor in the narrow sense,
it provides a useful non-neural structural ranking reference for identifying important or relevant papers.
Similarly, local heuristics such as Common Neighbors and Adamic--Adar
have long served as competitive, low-cost link prediction methods.
These approaches are especially valuable in a course project setting
because they are easy to implement, interpretable, and computationally inexpensive.
Including them helps us avoid an overly narrow comparison.
Instead, our evaluation can position GNNs against both embedding-based and heuristic alternatives,
which gives a more meaningful picture of what neural methods actually add.

For the specific task of link prediction with neural models,
one of the most relevant references is SEAL by Zhang and Chen \cite{zhang2018link}.
SEAL argues that link prediction should not rely solely on global node embeddings;
instead, it extracts the local enclosing subgraph around a candidate edge
and uses a GNN to classify whether that edge should exist.
This is conceptually important because it reframes link prediction
as a subgraph classification problem rather than a pure embedding reconstruction problem.
The paper reports strong improvements over latent feature methods such as
matrix factorization, node2vec, LINE, and variational graph autoencoders,
suggesting that the local subgraph structure contains predictive information
that simple embeddings may fail to capture.
SEAL is highly relevant to our proposal for two reasons.
First, it strengthens the case that GNNs are a suitable family of models for link prediction,
not only for node classification.
Second, it motivates our interest in structural encodings and positional information:
the success of SEAL depends heavily on making local structural roles visible to the model.
Although implementing full SEAL may be beyond the desired scope of a CPU-friendly course project,
its ideas justify our plan to experiment with simple structural or positional encodings
on top of GCN and GraphSAGE.

Another important thread in the literature concerns benchmarking and reproducibility.
Dwivedi et al. \cite{dwivedi2023benchmarking} argue that GNN research has often suffered
from inconsistent evaluation settings, non-comparable implementations,
and dataset choices that make it difficult to draw reliable conclusions.
Their benchmarking work emphasizes standardized pipelines,
careful separation of model classes, and reproducible experimentation.
This perspective is directly relevant to our project design.
Because we work within a course setting and limited compute budget,
it is important that our comparisons remain fair:
the same train/test splits, the same negative sampling strategy,
and comparable optimization settings should be used across methods whenever possible.
The benchmarking literature therefore does not merely provide background;
it shapes our methodology.
In particular, it supports our decision to use standard datasets,
fixed random seeds, and CPU-friendly baselines rather than overly complex architectures.

Finally, prior work on citation datasets themselves provides useful context
for interpreting model performance.
Sen et al. \cite{sen2008collective} and Namata et al. \cite{namata2012query}
document the characteristics of benchmark citation networks
and show that structure and attributes interact in nontrivial ways.
Their analyses suggest that citation graphs often have strong topical correlation along edges,
which helps explain why neighborhood aggregation methods perform well.
At the same time, they also highlight limitations of the datasets:
they are relatively small, sparse, and simplified compared to real scholarly recommendation systems.
This is an important caveat for our project.
Results on Cora or PubMed should be interpreted as evidence on controlled benchmarks,
not as definitive proof of effectiveness in recommending production-scale citations.

Overall, the literature suggests three main conclusions that motivate our project.
First, citation graphs are a well-established and meaningful domain for graph learning research.
Second, both classical methods and GNNs provide strong but conceptually different approaches
to link prediction.
Third, structural information beyond raw node embeddings --- such as local subgraphs,
positional encodings, or explicit directionality --- may substantially improve performance.
Our project builds on these insights by comparing classical heuristics,
node embedding methods, and simple GNN architectures in a unified and reproducible setting.
Rather than aiming for architectural novelty,
we focus on careful comparison and analysis of which components matter most
for citation link prediction under realistic course-level constraints.

\section{Methodology}

Our pipeline consists of four stages:
\begin{enumerate}
    \item \textbf{Graph preparation:} load the citation graph and construct both directed and symmetrized variants.
    \item \textbf{Edge split:} create reproducible train/validation/test splits for link prediction using a fixed random seed.
    \item \textbf{Model training:} train classical baselines and GNN encoders under comparable conditions.
    \item \textbf{Evaluation:} compare methods using classification- and ranking-based metrics on held-out edges.
\end{enumerate}

For the GNN models, we use an encoder--decoder setup.
The encoder computes node embeddings $z_u$ and $z_v$ from the node features and graph structure,
while the decoder assigns a link score to a candidate edge $(u,v)$.
As a simple baseline decoder, we will use either dot-product scoring or a small MLP over concatenated node embeddings.

To keep the comparison fair, all methods will use the same edge split, the same random seeds,
and the same negative sampling policy whenever possible.

\section{Baseline Methods}

We use three types of non-neural baselines:
\begin{itemize}
\item local heuristic methods such as Common Neighbors and Adamic--Adar,
\item node2vec embeddings with an edge scoring function,
\item an auxiliary PageRank-based ranking signal.
\end{itemize}

The first two categories serve as the main non-neural baselines against the GNN approaches,
while PageRank is included as a lightweight structural reference rather than as a direct pairwise link predictor.

\section{GNN Models}

We use two simple and well-established GNN encoders:

\begin{itemize}
\item \textbf{GCN:} a lightweight transductive message-passing baseline \cite{kipf2017semi}
\item \textbf{GraphSAGE:} an inductive neighborhood aggregation baseline \cite{hamilton2017inductive}
\end{itemize}

These models are deliberately simple.
The goal of the project is not to propose a new architecture,
but to determine whether standard GNN encoders already provide advantages on citation link prediction.

Both models will first be evaluated in a plain configuration.
Later, ablations will investigate whether their performance changes under directed graph modeling
and under additional structural or positional information.

\section{Research Extensions}

To keep the project focused, we define one primary and one secondary research lever.

\subsection{Primary Lever: Directed vs.\ Undirected Citation Modeling}
Citation graphs are naturally directional, but many GNN pipelines symmetrize the graph for convenience.
We will compare:
\begin{itemize}
    \item directed citation modeling,
    \item symmetrized (undirected) citation modeling.
\end{itemize}

\subsection{Secondary Lever: Structural or Positional Augmentation}
We will augment the baseline GNNs with lightweight structural or positional information,
for example, degree-based features or simple positional encodings.
The exact augmentation will be chosen so that it remains CPU-friendly and easy to reproduce.

\subsection{Planned Ablations}
\begin{itemize}
    \item plain GNN vs.\ GNN + structural/positional augmentation,
    \item directed vs.\ undirected graph construction,
    \item shallow vs.\ slightly deeper encoder configuration.
\end{itemize}

\section{Evaluation}

Because the original Planetoid benchmarks are primarily distributed with node-classification splits,
we will construct a fixed and reproducible \emph{edge-level} split for link prediction.
All compared methods will use the same train/validation/test partition.

For ranking-based evaluation, each source paper will be used to rank candidate target papers against sampled non-edges,
which allows us to report MRR and Hits@10 in addition to classification metrics.

\subsection{Primary Metrics}
\begin{itemize}
    \item ROC-AUC
    \item Average Precision (AP)
\end{itemize}

\subsection{Secondary Metrics}
\begin{itemize}
    \item MRR
    \item Hits@10
\end{itemize}

Accuracy, precision, recall, and F1 may also be reported for completeness,
but they will not be treated as the main metrics,
because link prediction depends strongly on the choice and ratio of negative samples.

\subsection{Fairness and Leakage Control}
Special care will be taken to avoid information leakage across splits,
especially in the undirected setting where reverse edges may otherwise appear in different partitions.

\section{Reproducibility Plan}

To support reproducibility, the project will use a GitLab repository,
fixed random seeds, a documented Python environment, and a shared PyTorch Geometric pipeline
for data loading, splitting, training, and evaluation.
All reported experiments will be designed to run on CPU under the course constraints.

\newpage
\section{Work Plan}

\begin{table}[ht]
\centering
\begin{tabular}{lp{10cm}}
\toprule
Week & Planned work \\
\midrule
5  & Finalize proposal, repository scaffold, dataset audit, and evaluation protocol. \\
6  & Implement data loading, reproducible edge splits, and heuristic baselines. \\
7  & Implement node2vec baseline and first evaluation pipeline. \\
8  & Implement GCN and GraphSAGE link prediction baselines. \\
9  & Run baseline experiments and first ablations on directionality. \\
10--11 & Add structural/positional augmentation and analyze secondary ablations. \\
12 & Consolidate results, tables, and reproducibility scripts. \\
13 & Final report, trained models/demo, and presentation. \\
\bottomrule
\end{tabular}
\caption{Planned timeline for the project.}
\end{table}

\section{Risks and Limitations}

\begin{itemize}
    \item \textbf{Benchmark simplicity:} Cora and PubMed are useful, but small, and may not reflect realistic recommendation settings.
    \item \textbf{Data leakage:} careless splitting can leak reverse-edge information between partitions.
    \item \textbf{Overfitting:} small graphs may favor overly tuned models.
    \item \textbf{Scope creep:} too many levers would weaken the experimental focus.
\end{itemize}

To mitigate these risks, we limit ourselves to simple baselines, fixed splits, small ablation sets,
and a single primary research lever.

\section{Conclusion}

This proposal defines a focused and reproducible study of link prediction in citation networks
using benchmark datasets from the Planetoid family.
Rather than introducing a novel architecture, the project aims to provide a careful comparison
between local heuristic methods, embedding-based baselines, and simple Graph Neural Networks.

The central research question is whether feature-aware GNN models such as GCN and GraphSAGE
offer clear advantages over classical approaches for citation link prediction,
and whether these results depend on modeling choices such as citation directionality
and lightweight structural or positional augmentation.

By using fixed edge-level splits, CPU-friendly models, and controlled ablations,
the project is designed to be feasible within the constraints of the course
while still producing a meaningful experimental analysis.
The expected outcome is a clear understanding of which components matter most
for citation link prediction on standard benchmark graphs.

\bibliographystyle{plain}
\bibliography{references}

\end{document}
